For each source and cumulant order define
\[
 \lambda_{j,\nu}=\operatorname{cum}_\nu(\varepsilon_j).
\]
The declared moment-determinacy assumption refers to the full moment sequence, so all moments and cumulants used below exist. The joint-cumulant partition identity is
\[
 \operatorname{cum}(Z_1,\ldots,Z_\nu)
 =\sum_{\pi\in\mathfrak P_\nu}
 (|\pi|-1)!(-1)^{|\pi|-1}
 \prod_{B\in\pi}\mathbb E\!\left[\prod_{s\in B}Z_s\right]. \tag{20}
\]
This is a finite multilinear polynomial. If the arguments split into two nonempty independent subfamilies, factorization of every moment in (20), followed by grouping partitions by the induced partitions on the two subfamilies, makes the alternating sum zero. Substituting \(X_i=\sum_j C_{ij}\varepsilon_j\), using multilinearity, and invoking \(\mathrm{ass:independent\text{-}sources}\) therefore gives the observable tensor identity
\[
 \mathcal T_\nu:=\operatorname{cum}_\nu(X)
 =\sum_{j=1}^n\lambda_{j,\nu}c_j^{\otimes\nu}
 =\sum_{j\in A_\nu}\lambda_{j,\nu}c_j^{\otimes\nu}. \tag{21}
\]

Fix \(\nu\ge2n-1\), put \(s=n-1\) and \(t=\nu-2n+2\), and reshape (21) into three blocks. Since \(n\ge p\ge2\), one has \(s\ge1\), and the order restriction gives \(t\ge1\). Thus
\[
 \widetilde{\mathcal T}_\nu
 =\sum_{j\in A_\nu}\lambda_{j,\nu}
 a_j\otimes b_j\otimes d_j,
 \qquad
 a_j=b_j=\operatorname{vec}(c_j^{\otimes s}),
 \qquad
 d_j=\operatorname{vec}(c_j^{\otimes t}). \tag{22}
\]
Let \(R=R_\nu\). Suppose first that \(R\ge2\). Because the columns of \(C\) are pairwise nonproportional and \(s=n-1\ge R-1\), \(\mathrm{lem:veronese\text{-}lagrange\text{-}rank}\) gives Kruskal rank \(R\) to each of the first two factor matrices in (22). Every pair of third-mode columns is linearly independent. Indeed, if \(z^{\otimes t}\) and \(w^{\otimes t}\) were proportional, then contraction against \(f^{\otimes t}\), for any linear functional \(f\) vanishing on \(w\), would give \(f(z)^t=0\). Every functional vanishing on \(w\) would then vanish on \(z\), implying \(z\in\operatorname{span}(w)\), contrary to distinct directions. Hence the third factor matrix has Kruskal rank at least two, and
\[
 k_1+k_2+k_3\ge R+R+2=2R+2. \tag{23}
\]
By \(\mathrm{lem:kruskal\text{-}three\text{-}factor\text{-}uniqueness}\), the decomposition (22) is unique up to a common permutation and reciprocal nonzero scalings.

The length \(R\) is minimal. Write \(A=[a_j]_{j\in A_\nu}\), \(B=[b_j]_{j\in A_\nu}\), and \(K=[b_j\otimes d_j]_{j\in A_\nu}\). The first two matrices have full column rank. If \(\sum_j q_jb_j\otimes d_j=0\), apply the row of a left inverse of \(B\) selecting column \(j\) to the first tensor component; this yields \(q_jd_j=0\), and \(d_j\ne0\) yields \(q_j=0\). Thus \(K\) also has full column rank. The first matricization of (22) is
\[
 A\operatorname{diag}(\lambda_{j,\nu})K^\top, \tag{24}
\]
whose matrix rank is \(R\), because both outer factors have column rank \(R\) and every \(\lambda_{j,\nu}\), \(j\in A_\nu\), is nonzero by definition. A sum of fewer than \(R\) rank-one tensors has first-matricization rank below \(R\), contradicting (24). Thus the minimum number of terms is exactly \(R\), and Kruskal uniqueness implies that every least feasible decomposition in \(\mathrm{def:multi\text{-}order\text{-}handle}\) returns precisely \(\{[c_j]:j\in A_\nu\}\).

This minimization is exposed as a finite population-oracle operation rather than an opaque factorization. At fixed \(\nu\) and \(R\), its variables are \(\gamma_1,\ldots,\gamma_R\in\mathbb R\) and \(z_1,\ldots,z_R\in\mathbb R^p\). Its constraints are the finitely many coordinate equations
\[
 \widetilde{\mathcal T}_\nu
 =\sum_{h=1}^{R}\gamma_h
 \operatorname{vec}(z_h^{\otimes(n-1)})\otimes
 \operatorname{vec}(z_h^{\otimes(n-1)})\otimes
 \operatorname{vec}(z_h^{\otimes(\nu-2n+2)}), \tag{25}
\]
together with the polynomial inequalities \(\gamma_h^2>0\) and \(\lVert z_h\rVert_2^2>0\). The handle asks its exact feasibility oracle about (25) for the finite list \(R=0,\ldots,n\) and chooses the least feasible value. Exact projective equality of two returned nonzero vectors is tested by the vanishing of all two-by-two minors of the matrix having those vectors as columns. Thus every operation at a processed order is explicit. The population law supplies exact cumulants; no claim is made that these oracle operations constitute a finite-sample numerical algorithm.

If \(R=0\), (21) is the zero tensor and the empty decomposition is the unique least feasible output. If \(R=1\), it is a nonzero rank-one tensor. The projective Veronese map is injective: the contraction argument used after (22) shows that proportional nonzero tensor powers imply proportional base vectors. Hence the unique active direction is recovered. This covers all active ranks.

The same proof gives the adaptive ex-post certificate. If \(R_\nu\ge2\) and \(\nu\ge2R_\nu-1\), use degrees
\[
 (R_\nu-1,R_\nu-1,\nu-2R_\nu+2). \tag{26}
\]
The first two degrees are at least \(R_\nu-1\), and the third is positive. The first two Kruskal ranks are therefore \(R_\nu\), the third is at least two, and (23)--(24) prove uniqueness and minimality. Because \(R_\nu\) is learned only after solving the order-\(\nu\) systems, (26) is an ex-post certificate. The universally executable schedule in \(\mathrm{def:multi\text{-}order\text{-}handle}\) uses the known \(n\) and begins at \(2n-1\); it does not presuppose \(R_\nu\).

By \(\mathrm{lem:marcinkiewicz\text{-}cumulant\text{-}tail}\), every declared nondegenerate non-Gaussian source has infinitely many nonzero cumulants of order at least three. This set of orders is unbounded, so for each \(j\) there is a finite \(\nu_j\ge2n-1\) with \(\lambda_{j,\nu_j}\ne0\). Consequently
\[
 K_{C,\varepsilon}=\max_{1\le j\le n}\nu_j<\infty. \tag{27}
\]
By the orderwise recovery proved above, processing \(2n-1,2n,\ldots,K_{C,\varepsilon}\) and unioning the outputs includes every \([c_j]\). Pairwise nonproportionality prevents erroneous projective deduplication, and no direction outside this set is returned. The handle therefore stops no later than (27), when the union first has cardinality \(n\), and returns
\[
 \{[c_1],\ldots,[c_n]\}=\mathcal D(P_X), \tag{28}
\]
where the equality uses \(\mathrm{thm:sharp\text{-}effect\text{-}frontier}\). No cross-order labeling is needed. Choosing an arbitrary nonzero representative of each recovered direction produces \(C'=C\Pi D\); \(\mathrm{lem:frontier\text{-}monomial\text{-}invariance}\) then yields the same \(\mathcal V_x\) and \(\mathcal R_{xy}\). The sharp-effect theorem, rather than a definition, proves that these evaluated sets equal the causal law-fiber response targets.

We next prove that no cumulant cutoff is uniform on the qualitative class. Fix any finite \(K\ge3\) and \(\rho>0\). By \(\mathrm{lem:finite\text{-}moment\text{-}near\text{-}gaussian\text{-}witness}\), there is a centered, variance-one, non-Gaussian, moment-determinate law \(F\) within total variation \(\rho\) of \(G=\mathcal N(0,1)\), whose moments through order \(K\) agree with those of \(G\). Give two sources independent law \(F\), fix \(0<\gamma<\pi/4\), and set
\[
 C_0=I_2,
 \qquad
 C_1=
 \begin{pmatrix}
 \cos\gamma&-\sin\gamma\\
 \sin\gamma&\cos\gamma
 \end{pmatrix}. \tag{29}
\]
Both matrices are normalized and full rank and have nonzero pairwise nonproportional columns. Independence makes all source joint moments of total degree at most \(K\) equal the corresponding moments under \(G^{\otimes2}\). Since \(C_1\) is orthogonal and \(G^{\otimes2}\) is orthogonally invariant, the two observed laws in (29) have identical moments, hence identical cumulants, through order \(K\), while their projective direction sets differ. Therefore no fixed finite cumulant cutoff recovers the directions throughout the qualitative class.

For the statistical boundary, for each \(N\) choose the preceding witness with
\[
 d_{\mathrm{TV}}(F_N,G)\le\frac{1}{4N^2}, \tag{30}
\]
and let \(P_{k,N}\) be the law generated by \(C_k\) and two independent \(F_N\) sources. Total variation contracts under measurable pushforward because the preimage of each measurable event is measurable. Also, a telescoping expansion of product signed measures and the triangle inequality give
\[
 d_{\mathrm{TV}}(\mu^{\otimes m},\nu^{\otimes m})
 \le m\,d_{\mathrm{TV}}(\mu,\nu). \tag{31}
\]
Each observed law is within \(2d_{\mathrm{TV}}(F_N,G)\) of the common law \(G^{\otimes2}\), the latter remaining unchanged by either orthogonal matrix in (29). The triangle inequality, (30), and then (31) imply
\[
 d_{\mathrm{TV}}(P_{0,N},P_{1,N})\le N^{-2},
 \qquad
 d_{\mathrm{TV}}(P_{0,N}^{\otimes N},P_{1,N}^{\otimes N})\le N^{-1}\longrightarrow0. \tag{32}
\]

The two projective direction sets in (29) have a fixed positive Hausdorff separation. A uniformly consistent direction-set estimator would yield a test by selecting the closer of these two fixed targets, with both errors tending to zero. For every test event \(E\), however, the definition of total variation gives
\[
 P_{0,N}^{\otimes N}(E^c)+P_{1,N}^{\otimes N}(E)
 \ge1-d_{\mathrm{TV}}(P_{0,N}^{\otimes N},P_{1,N}^{\otimes N}), \tag{33}
\]
which tends to one by (32), a contradiction.

The corresponding causal frontiers are fixed and separated. Directly applying \(\mathrm{def:admissible\text{-}source\text{-}set}\) at \(x=1\), \(y=2\), gives
\[
 \mathcal R_{12}(C_0)=\{0\},
 \qquad
 \mathcal R_{12}(C_1)=\{\tan\gamma,-\cot\gamma\}. \tag{34}
\]
Thus the nearest-target argument also rules out a uniformly consistent frontier estimator. More strongly, let \(\mathcal U_N\) be a random compact set with uniform frontier coverage at least \(1-\beta\), where \(\beta<1/2\). Under \(P_{0,N}^{\otimes N}\), coverage of the first target has probability at least \(1-\beta\). Coverage of the second target has probability at least \(1-\beta\) under \(P_{1,N}^{\otimes N}\), and (32) transfers that event to \(P_{0,N}^{\otimes N}\) with probability at least \(1-\beta-o(1)\). Their intersection consequently has probability at least \(1-2\beta-o(1)>0\). On the intersection, since \(0\in\mathcal R_{12}(C_0)\) and \(-\cot\gamma\in\mathcal U_N\),
\[
 d_H(\mathcal U_N,\mathcal R_{12}(C_0))\ge\cot\gamma. \tag{35}
\]
Hence no uniformly honest frontier-containing confidence union with coverage strictly above one half has uniformly vanishing Hausdorff error.

Finally, two explicit families establish the deterministic discontinuity boundaries. Column normalization may be applied because \(\mathrm{lem:frontier\text{-}monomial\text{-}invariance}\) preserves the frontiers. For
\[
 D_t=\begin{pmatrix}t&1\\1&1\end{pmatrix},
\]
the matrices are full rank and have nonzero nonproportional columns for all sufficiently small \(|t|\) with \(t\ne1\), and deletion-rank calculation gives
\[
 \mathcal R_{12}(D_t)=\{1/t,1\}\quad(t\ne0),
 \qquad
 \mathcal R_{12}(D_0)=\{1\}. \tag{36}
\]
This proves discontinuity at a vanishing treatment loading. For
\[
 E_t=
 \begin{pmatrix}
 1&1&0&1\\
 1&t&0&0\\
 0&0&1&1
 \end{pmatrix},
\]
the matrices are full row rank with nonzero nonproportional columns for all sufficiently small \(|t|\). Deleting row one and the indicated column shows
\[
 \mathcal R_{12}(E_0)=\{0\},
 \qquad
 \mathcal R_{12}(E_t)=\{0,t,1\}\quad(t\ne0). \tag{37}
\]
Indeed, at \(t=0\) only columns two and four are admissible and both ratios are zero; at \(t\ne0\), columns one, two, and four are admissible and give ratios \(1,t,0\). This proves discontinuity at a deletion-rank boundary.

Equations (32), (36), and (37) show that the qualitative assumptions alone supply neither a uniform stopping order nor a uniform statistical modulus. Quantitative non-Gaussianity, moment, Veronese-conditioning, angular, probe/eigen-gap, denominator, and deletion-rank margins therefore cannot be omitted from a uniform root-\(N\) result. On the declared separated common-order iid class, \(\mathrm{thm:uniform\text{-}confidence\text{-}frontier}\) supplies simultaneous response-vector and scalar coverage with \(O(N^{-1/2})\) Hausdorff error. This completes the pointwise population-oracle construction and the nonuniform inference boundary without strengthening any assumption.
